\documentclass{article}
\usepackage[utf8]{inputenc}
\usepackage[margin=1in]{geometry}
\usepackage{hyperref}
\usepackage{listings}
\usepackage{graphicx}
\usepackage{booktabs}
\usepackage{longtable}
\usepackage{enumitem}
\usepackage{tikz}
\usepackage{pgf-umlsd}
\usepgflibrary{arrows}

\hypersetup{
  colorlinks=true,
  linkcolor=blue,
  urlcolor=blue,
  citecolor=blue
}

\lstset{
  basicstyle=\ttfamily\small,
  breaklines=true,
  frame=single,
  numbers=left,
  numberstyle=\tiny,
  showstringspaces=false
}

\title{\textbf{Technical Design Document}\\[0.5em]
\large IntraCore Solutions\\
\large Voice Assistant for IT Ticketing and Identity Verification\\
\large Version 1.0.0\\
\large March 5, 2026}
\author{}
\date{}

\begin{document}

\maketitle

\begin{abstract}
This document describes the technical design of a voice assistant for IT ticketing and identity verification, built for IntraCore Solutions as a Cognigy Capstone project. The system combines Cognigy.AI conversational AI with xApps for visual interaction, backed by a REST API that integrates with PostgreSQL for persistent data. Users can verify their identity via voice (ANI/phone, email, employee number), manage internal support tickets, and access FAQ content through WebRTC or dial-in voice channels in multiple languages.
\end{abstract}

\tableofcontents
\newpage

\section{Introduction}

\subsection{Purpose}
This Technical Design Document (TDD) provides a structured walkthrough of the Cognigy Capstone solution for IntraCore Solutions. It serves as the primary reference for understanding the architecture, APIs, data model, and deployment of the voice assistant system.

\subsection{Scope}
The document covers the full technical implementation: REST API backend, Cognigy voice agent integration, xApp interfaces, identity verification, ticketing workflows, and deployment infrastructure.

\subsection{Audience}
This document is intended for capstone evaluators, customer stakeholders, and technical team members who need to understand or extend the system.

\section{Business Context}

IntraCore Solutions is a fictional mid-sized technology firm with distributed teams across multiple locations. The organization operates a shared internal support center that handles IT requests, facilities issues, and general inquiries. The voice assistant project aims to:

\begin{itemize}[noitemsep]
  \item Automate common support interactions via voice
  \item Provide secure identity verification before sensitive operations
  \item Support multilingual employees (English, Thai, Japanese, German)
  \item Reduce load on human agents for routine tasks
\end{itemize}

\section{Capstone Scope}

The implemented scope includes:

\begin{itemize}[noitemsep]
  \item \textbf{Identity \& Verification (ID\&V):} Verification via ANI/phone number, email, and employee number; matches against the employees table
  \item \textbf{Internal Ticketing:} Open (create), Check Status, Update, Close, and Search History
  \item \textbf{FAQ:} Common questions answered by the voice agent
  \item \textbf{Voice Channels:} WebRTC and dial-in
  \item \textbf{Multilingual:} Support for multiple languages in both voice and xApps
\end{itemize}

\section{System Overview}

\subsection{Architecture}
The system follows a three-tier architecture:

\begin{enumerate}[noitemsep]
  \item \textbf{Cognigy Voice Agent} --- Handles speech recognition, NLU, and dialog management; invokes REST API for backend operations
  \item \textbf{REST API} --- Elysia-based backend on Bun; validates requests, enforces auth, and performs business logic
  \item \textbf{PostgreSQL} --- Persistent storage for companies, employees, tickets, sessions, and whitelists
\end{enumerate}

\subsection{Flow}
\begin{enumerate}[noitemsep]
  \item User connects via WebRTC or dial-in
  \item Cognigy agent authenticates via GitHub OAuth $\rightarrow$ POST /v1/auth/token
  \item Agent calls POST /v1/auth/verify with phone, email, employee number for ID\&V
  \item Agent calls ticket APIs (create, list, update, close, history) with Bearer token
  \item xApps (verification, ticket-list, ticket-detail) are shown in Cognigy when needed; they use the app-page SDK to submit data back to the agent
\end{enumerate}

\subsection{Sequence Diagrams}

The following diagrams depict the flow (user connects via a public phone number or WebRTC).

\subsubsection{Authentication Flow}
\begin{figure}[htbp]
  \centering
  \begin{sequencediagram}
    \newthread{user}{User}
    \newinst{agent}{Cognigy}
    \newinst{api}{IntraCore API}
    \newinst{gh}{GitHub}
    \newinst{db}{Database}
    \begin{call}[1]{user}{Make call}{agent}{Authenticated}
      \begin{call}[2]{agent}{POST /v1/auth/token}{api}{accessToken, expiresAt}
        \begin{call}[1]{api}{GET user}{gh}{email, profile}
        \end{call}
        \begin{call}[1]{api}{SELECT whitelists}{db}{allowed?}
        \end{call}
        \begin{call}[1]{api}{INSERT sessions}{db}{OK}
        \end{call}
      \end{call}
    \end{call}
  \end{sequencediagram}
  \caption{Authentication: User makes call to Cognigy; Cognigy calls POST /v1/auth/token; API validates with GitHub, checks whitelist, creates session; Cognigy proceeds with authenticated user. Note: githubToken is supplied by the Cognigy actor.}
\end{figure}

\subsubsection{Identity and Verification Flow}
\begin{figure}[htbp]
  \centering
  \begin{sequencediagram}
    \newthread{user}{User}
    \newinst{xapp}{xApp}
    \newthread{agent}{Cognigy}
    \newinst{api}{IntraCore API}
    \newinst{db}{Database}

    \mess{agent}{Email xApp link}{user}
    \begin{call}[1]{user}{Open xApp link}{user}{}
    \end{call}
    \begin{call}[1]{user}{Submit credentials}{xapp}{}
      \begin{call}[1]{xapp}{Forward credentials}{agent}{}
      \end{call}
    \end{call}
    \begin{call}[2]{agent}{POST /v1/auth/verify}{api}{id, fullName, preferredLanguage}
      \begin{call}[1]{api}{Validate session}{db}{valid?}
      \end{call}
      \begin{call}[1]{api}{Match employee}{db}{employee?}
      \end{call}
    \end{call}
  \end{sequencediagram}
  \caption{ID\&V: User requests identity verification; Cognigy emails xApp link; user opens link, submits credentials to xApp; xApp forwards to Cognigy; API matches against employees table; Cognigy returns verified identity to user.}
\end{figure}

\subsubsection{Ticket Creation Flow}
\begin{figure}[htbp]
  \centering
  \begin{sequencediagram}
    \newthread{user}{User}
    \newinst{agent}{Cognigy}
    \newinst{api}{IntraCore API}
    \newinst{db}{Database}
    \begin{call}[1]{user}{Create ticket}{agent}{ticket}
      \begin{call}[1]{agent}{Prompt: describe issue}{user}{OK}
      \end{call}
      \begin{call}[2]{agent}{POST /v1/tickets}{api}{ticket}
        \begin{call}[1]{api}{Validate company, reporter}{db}{OK}
        \end{call}
        \begin{call}[1]{api}{Count tickets}{db}{count}
        \end{call}
        \begin{call}[1]{api}{Insert ticket}{db}{OK}
        \end{call}
      \end{call}
    \end{call}
  \end{sequencediagram}
  \caption{Ticket creation: User requests ticket creation; Cognigy prompts for issue description; API validates and inserts ticket; Cognigy returns ticket to user.}
\end{figure}

\section{Technology Stack}

\begin{itemize}[noitemsep]
  \item \textbf{Runtime:} Bun
  \item \textbf{Web Framework:} Elysia
  \item \textbf{ORM:} Drizzle
  \item \textbf{Database:} PostgreSQL
  \item \textbf{Logging:} Pino
  \item \textbf{xApp UI:} Material Web (md-* components)
\end{itemize}

\section{Data Model}

\subsection{companies}
\begin{lstlisting}[language=SQL]
id          TEXT PRIMARY KEY
slug        TEXT NOT NULL UNIQUE
name        TEXT NOT NULL
description TEXT
created_at  TEXT NOT NULL
updated_at  TEXT NOT NULL
deleted_at  TEXT
\end{lstlisting}

\subsection{employees}
\begin{lstlisting}[language=SQL]
id                 TEXT PRIMARY KEY
employee_number    TEXT NOT NULL
full_name          TEXT NOT NULL
email              TEXT NOT NULL UNIQUE
phone_number       TEXT NOT NULL
department         TEXT
role               TEXT
preferred_language TEXT NOT NULL DEFAULT 'en-US'
company_id         TEXT NOT NULL REFERENCES companies(id)
created_at         TEXT NOT NULL
updated_at         TEXT NOT NULL
deleted_at         TEXT
UNIQUE(company_id, employee_number)
\end{lstlisting}

\subsection{tickets}
\begin{lstlisting}[language=SQL]
id             TEXT PRIMARY KEY
ticket_number  TEXT NOT NULL UNIQUE
title          TEXT NOT NULL
description    TEXT
status         TEXT NOT NULL DEFAULT 'NEW'
  -- NEW | PENDING | RESOLVED | CANCELLED | CLOSED
priority       TEXT NOT NULL DEFAULT 'MEDIUM'
  -- LOW | MEDIUM | HIGH
category       TEXT
  -- IT | FACILITIES | MISCELLANEOUS
assignee_id    TEXT REFERENCES employees(id)
reported_by_id TEXT NOT NULL REFERENCES employees(id)
company_id     TEXT NOT NULL REFERENCES companies(id)
created_at     TEXT NOT NULL
updated_at     TEXT NOT NULL
closed_at      TEXT
\end{lstlisting}

\subsection{sessions}
\begin{lstlisting}[language=SQL]
id          TEXT PRIMARY KEY
token_hash  TEXT NOT NULL UNIQUE
email       TEXT NOT NULL
expires_at  TEXT NOT NULL
created_at  TEXT NOT NULL
\end{lstlisting}

\subsection{whitelists}
\begin{lstlisting}[language=SQL]
id         TEXT PRIMARY KEY
email      TEXT NOT NULL UNIQUE
created_at TEXT NOT NULL
\end{lstlisting}

\section{API Design}

\subsection{Public Endpoints}
\begin{itemize}[noitemsep]
  \item \texttt{GET /} --- Health/root response
  \item \texttt{GET /health} --- Health check
  \item \texttt{POST /v1/auth/token} --- Exchange GitHub token for session (body or Bearer)
\end{itemize}

\subsection{Protected Endpoints (Bearer Auth)}
\begin{itemize}[noitemsep]
  \item \texttt{POST /v1/auth/verify} --- Identity verification (phoneNumber, email, employeeNumber)
  \item \texttt{GET /v1/tickets} --- List tickets
  \item \texttt{POST /v1/tickets} --- Create ticket
  \item \texttt{GET /v1/tickets/:id} --- Get ticket by ID
  \item \texttt{PUT /v1/tickets/:id} --- Update ticket
  \item \texttt{DELETE /v1/tickets/:id} --- Close ticket
  \item \texttt{GET /v1/tickets/number/:ticketNumber} --- Get ticket by ticket number
  \item \texttt{PUT /v1/tickets/number/:ticketNumber} --- Update ticket by number
  \item \texttt{DELETE /v1/tickets/number/:ticketNumber} --- Close ticket by number
  \item \texttt{GET /v1/tickets/history?employeeId=} --- Ticket history for employee
  \item \texttt{GET|POST /v1/companies}, \texttt{GET|PUT|DELETE /v1/companies/:id}
  \item \texttt{GET|POST /v1/employees}, \texttt{GET|PUT|DELETE /v1/employees/:id}
  \item \texttt{GET|POST /v1/whitelists}, \texttt{GET|PUT|DELETE /v1/whitelists/:id}
  \item \texttt{GET /v1/sessions}, \texttt{GET|DELETE /v1/sessions/:id}
\end{itemize}

\subsection{Use Case to API Mapping}
\begin{longtable}{@{}ll@{}}
\toprule
\textbf{Use Case} & \textbf{API} \\
\midrule
\endfirsthead

\toprule
\textbf{Use Case} & \textbf{API} \\
\midrule
\endhead

ID\&V & POST /v1/auth/verify \\
Open ticket & POST /v1/tickets \\
Check status & GET /v1/tickets/number/:ticketNumber or GET /v1/tickets/:id \\
Update ticket & PUT /v1/tickets/:id or PUT /v1/tickets/number/:ticketNumber \\
Close ticket & DELETE /v1/tickets/:id or DELETE /v1/tickets/number/:ticketNumber \\
Search history & GET /v1/tickets/history?employeeId= \\
\bottomrule
\end{longtable}

\subsection{OpenAPI}
Interactive API documentation is available at \texttt{/docs}.

\section{Identity and Verification}

\texttt{POST /v1/auth/verify} requires a valid Bearer token and accepts:

\begin{lstlisting}
{
  "phoneNumber": "+1234567890",
  "email": "user@example.com",
  "employeeNumber": "12345"
}
\end{lstlisting}

The API matches all three fields against the \texttt{employees} table (excluding soft-deleted rows). On success it returns:

\begin{lstlisting}
{
  "id": "uuid",
  "fullName": "Jane Doe",
  "email": "user@example.com",
  "preferredLanguage": "en-US"
}
\end{lstlisting}

On failure: 401 with \texttt{\{"error":"Unauthorized","message":"Identity verification failed"\}}.

\section{Authentication}

\begin{enumerate}[noitemsep]
  \item User authenticates with GitHub OAuth (outside this API)
  \item Client sends GitHub token to \texttt{POST /v1/auth/token} (in body as \texttt{githubToken} or in \texttt{Authorization: Bearer <token>})
  \item API validates GitHub token and checks email against \texttt{whitelists}
  \item If whitelisted, API creates a session and returns \texttt{accessToken}, \texttt{expiresAt}, \texttt{tokenType: "Bearer"}
  \item All protected routes require \texttt{Authorization: Bearer <accessToken>}
\end{enumerate}

\section{Cognigy xApps}

Three xApps are provided, built with the Cognigy app-page SDK and Material Web:

\begin{itemize}[noitemsep]
  \item \textbf{verification.html} --- Collects phone number, email, employee number; submits via \texttt{SDK.submit(\{phoneNumber, email, employeeNumber\})}
  \item \textbf{ticket-list.html} --- Displays ticket history with status/priority badges; supports en, th, ja, de
  \item \textbf{ticket-detail.html} --- Shows full ticket details (status, priority, category, description, company, reporter, assignee, dates)
\end{itemize}

All xApps use \texttt{/sdk/app-page-sdk.js}, Material Web components (\texttt{md-outlined-text-field}, \texttt{md-filled-button}, \texttt{md-outlined-select}), and multilingual labels.

\section{Deployment}

\subsection{Docker}
Multi-stage build:

\begin{lstlisting}[language=bash]
# Stage 1: Build
FROM oven/bun:1 AS builder
WORKDIR /app
COPY package.json bun.lock ./
RUN bun install --frozen-lockfile
COPY . .
RUN bun build --compile --minify server.ts --outfile intracore

# Stage 2: Runtime
FROM gcr.io/distroless/cc-debian12:nonroot
WORKDIR /app
COPY --from=builder --chown=65532:65532 /app/intracore .
EXPOSE 3000
ENV PORT=3000
CMD ["/app/intracore"]
\end{lstlisting}

\subsection{Environment}
\begin{itemize}[noitemsep]
  \item \texttt{DATABASE\_URL} --- PostgreSQL connection string
  \item \texttt{PORT} --- Server port (default 3000)
\end{itemize}

\section{CI/CD}

GitHub Actions workflow \texttt{.github/workflows/docker-publish.yml}:

\begin{itemize}[noitemsep]
  \item Triggers on push to \texttt{main}
  \item Logs in to \texttt{ghcr.io}
  \item Builds and pushes image to \texttt{ghcr.io/boss-spicyz100x/intracore:latest} and \texttt{ghcr.io/boss-spicyz100x/intracore:<sha>}
\end{itemize}

\section{Security}

\begin{itemize}[noitemsep]
  \item \textbf{Token hashing:} Session tokens are hashed (SHA-256) before storage; only the hash is stored
  \item \textbf{Whitelist:} Only whitelisted emails can obtain session tokens
  \item \textbf{HTTPS:} Production deployments should use HTTPS (handled by ingress/load balancer)
\end{itemize}

\section{Deliverables Summary}

\begin{itemize}[noitemsep]
  \item Voice agent with Cognigy flow design
  \item REST API integration for ID\&V and ticketing
  \item xApp (verification, ticket-list, ticket-detail)
  \item Multilingual support (en, th, ja, de)
  \item Presentations and documentation for capstone evaluation
\end{itemize}

\end{document}
